\chapter[Terminal, naturalité et limite]{Terminal multivalué, naturalité unique et passage à la limite}
\label{chap:pdfv2-terminal-naturalidad-limite}

\begin{thesisbox}
Le continu discret produit un seul système d'histoires, d'extensions, de mémoires et de terminaux. Les discriminants internes \(\mathrm{sp},\mathrm{sol},\mathrm{gau},\mathrm{coh},\mathrm{det}\) désignent cinq caractéristiques simultanées de ce même système : ils ne fixent pas d'entrées distinctes et n'autorisent pas à choisir des histoires différentes. Dans ce chapitre, le terminal conserve toutes les continuations, la connexion nonadique reste une correspondance et la naturalité s'exprime par une seule identité.
\end{thesisbox}

\section{Histoires survivantes avec témoins d'extension}

Soit \(\widehat x_k\in\widehat{\mathcal X}^{\rm tot}_k\) l'état total de \eqref{eq:continuo-comun-dominio-total}. Pour \(N\geq k\), nous définissons sans autre prédicat d'appartenance
\begin{equation}
\label{eq:pdfv2-terminal-historias-finita}
 \mathscr H_{k,N}(\widehat x_k)
 :=\operatorname{Term}_{k,N}(\widehat x_k).
\end{equation}
Chaque \(h=(\varepsilon_{k+1},\ldots,\varepsilon_N)\) du membre de droite contient, par \eqref{eq:continuo-comun-vertice-arbol}, ses états intermédiaires, les témoins de \(\operatorname{Ext}\), les retenues, les frontières, les parcours et le registre. Chaque arête satisfait l'action de phase \eqref{eq:continuo-comun-accion-fase-tpk} ; tout sous-chemin de neuf pas détermine donc, avec ses neuf émissions et non seulement sa longueur, un élément de \(\mathsf Z^\Gamma_{9,j}\) et un témoin de \(\Gamma_{9,j}\) par \eqref{eq:continuo-comun-gamma9-apice}--\eqref{eq:continuo-comun-gamma9-relacion}. Ainsi, Ext et l'holonomie nonadique naissent du même mot compatible avec la phase ; on n'ajoute à \eqref{eq:pdfv2-terminal-historias-finita} ni filtre de survie ni seconde condition \(\Gamma_9\).

Pour \(k\leq\ell\leq N\), soit
\begin{equation}
\label{eq:pdfv2-terminal-prefijos-intermedios}
 \mathscr P_{\ell\mid k}(\widehat x_k)
 :=\{h_{[k,\ell]}:h\in\mathscr H_{k,N}(\widehat x_k)
       \text{ pour un certain }N\geq\ell\}.
\end{equation}
Pour \(p\in\mathscr P_{\ell\mid k}(\widehat x_k)\), nous utilisons l'abréviation
\begin{equation}
\label{eq:pdfv2-terminal-continuaciones-prefijo}
 \mathscr H_{\ell,N}(p)
 :=\{q:p\star q\in\mathscr H_{k,N}(\widehat x_k)\}.
\end{equation}
Le prolongement neutre de \eqref{eq:continuo-comun-seccion-total} étend chaque préfixe à tout horizon. En conséquence,
\begin{equation}
\label{eq:pdfv2-terminal-no-vacuidad-historias}
 \mathscr H_{k,N}(\widehat x_k)\neq\varnothing,
 \qquad
 \mathscr P_{\ell\mid k}(\widehat x_k)\neq\varnothing,
\end{equation}
par \cref{prop:continuo-comun-arbol-finito-no-vacio}.

La troncature
\begin{equation}
\label{eq:pdfv2-terminal-truncamiento-historias}
 \pi_{N+1,N}^{\mathscr H}:
 \mathscr H_{k,N+1}(\widehat x_k)
 \longrightarrow\mathscr H_{k,N}(\widehat x_k)
\end{equation}
est exactement \eqref{eq:continuo-comun-proyeccion-terminal} et est surjectif par \cref{lem:continuo-comun-terminal-sobreyectivo}. La survie est, par \eqref{eq:continuo-comun-surv-igual-total}, une sortie démontrée du domaine total, non une sélection préalable.

\section{Registre terminal engendré et ses \(13\) coordonnées}

Fixons une histoire
\[
 h=(\varepsilon_{j+1}:x_j\longrightarrow x_{j+1})_{j=k}^{N-1},
 \qquad
 \gamma_{k,k}:=I_{x_k},\qquad
 \gamma_{k,j+1}:=\gamma_{k,j}\star\varepsilon_{j+1},
\]
La notation opératorielle \(U_{\gamma_{k,j+1}}
=U_{\varepsilon_{j+1}}U_{\gamma_{k,j}}\) conserve l'ordre d'application, tandis que \(\star\) conserve l'ordre chronologique du mot. Nous fixons en outre un niveau de coefficients \(r\geq1\). Pour abréger, nous écrivons
\[
 C_{x_j,r}:=C_{j,r}(x_j),\qquad
 F_{\varepsilon_{j+1}}:C_{x_{j+1},r}\longrightarrow C_{x_j,r}
\]
pour le complexe et l'application de complexes de \eqref{eq:pdfv2-cor-cono-celda} et de \eqref{eq:pdfv2-cor-refinamiento-celda}. Le type du registre de niveau \(j\), noté \(\mathfrak R_{j,r}\), est un registre dépendant possédant les treize champs suivants. Pour chaque champ, nous indiquons le type, la graine et l'action d'une arête \(\varepsilon:=\varepsilon_{j+1}\).

L'état \(x_k\) contient déjà le mot complet \(\gamma_{0,k}=\varepsilon_1\star\cdots\star\varepsilon_k\). Ainsi, une ``graine en \(k\)'' ne réinitialise jamais la généalogie : elle est le pliage des mêmes actualisations depuis la racine jusqu'à \(k\). Nous notons
\begin{equation}
\label{eq:pdfv2-terminal-prefijo-heredado}
 \begin{gathered}
 M_k^\pm:=\sum_{i=1}^{k}b_{\varepsilon_i}^\pm,qquad
 D_k:=\sum_{i=1}^{k}b_{\varepsilon_i},qquad
 H_k:=\sum_{i=1}^{k}|b_{\varepsilon_i}|,\\
 \mathbf c_{\leq k}:=
  \bigl(c_r(\varepsilon_i,\gamma_{0,i-1})\bigr)_{1\leq i\leq k},
 \qquad
 \mathbf k_{\leq k}:=
  \bigl(\kappa_{i-1}^{\rm can}\bigr)_{1\leq i\leq k},\\
 \mathbf F_{\leq k}:=
  \bigl(F_{\varepsilon_i}:C_{x_i,r}\to C_{x_{i-1},r}\bigr)_{1\leq i\leq k}.
 \end{gathered}
\end{equation}
Toutes les listes vides ci-dessous ne sont à entendre littéralement qu'en \(k=0\) ; pour \(k>0\), elles sont remplacées par les préfixes hérités de \eqref{eq:pdfv2-terminal-prefijo-heredado} et du registre de \(\gamma_{0,k}\). Ainsi, le registre ne conserve pas de copie opaque de \(x_k\), mais ne perd pas non plus son passé APP--TRIT--TPK.

\paragraph{Fibre \(\mathsf F\) : domaine local et conservation sous prolongement}
La fibre élémentaire est
\begin{equation}
\label{eq:pdfv2-terminal-tipo-fibra}
 \operatorname{Fib}_r(x):=
 \bigl(\mathsf H_x^\Sigma\oplus\mathsf H_x^\Pi,\,
       \mathsf M_x,\,
       \operatorname{Cell}_{\rm TPK}(x),\,C_{x,r}\bigr).
\end{equation}
Le champ \(\mathsf F_{k,j}\) est un mot dans la catégorie de ces fibres. Sa graine et son actualisation sont
\begin{align}
 \mathsf F_{k,k}
 &:=\bigl(\operatorname{Fib}_r(x_0),I\bigr)\star
 \mathop{\star}_{i=1}^{k}
 \bigl(\operatorname{Fib}_r(x_i),U_{\varepsilon_i},
 \mathsf C_{\varepsilon_i},\operatorname{Cell}_{\rm TPK}(\varepsilon_i),
 F_{\varepsilon_i}\bigr),
 \label{eq:pdfv2-terminal-semilla-f}\\
 \operatorname{Upd}^{\mathsf F}_{\varepsilon}\mathsf F_{k,j}
 &:=
 \mathsf F_{k,j}\star
 \bigl(\operatorname{Fib}_r(x_{j+1}),U_\varepsilon,
 \mathsf C_\varepsilon,\operatorname{Cell}_{\rm TPK}(\varepsilon),
 F_\varepsilon\bigr).
 \label{eq:pdfv2-terminal-update-f}
\end{align}
Elle est non vide puisque les deux feuillets, la mémoire, l'unité cellulaire \eqref{eq:pdfv2-cor-cell-unidad} et le complexe des ports sont déjà construits. La composition de ses cinq applications prouve la naturalité.

\paragraph{Extensions \(\mathsf E\) : morphismes de prolongement compatible de l'état}
Le codomaine est la catégorie des mots de relations typées. Nous posons
\begin{align}
 \mathsf E_{k,k}
 &:=(I_{x_0},\operatorname{Ext}_{I_{x_0}})\star
 \mathop{\star}_{i=1}^{k}
 \bigl(\varepsilon_i,\operatorname{Ext}_{\varepsilon_i},
 \operatorname{Ext}_{\varepsilon_i}\circ
 \operatorname{Ext}_{\gamma_{0,i-1}}
 \subseteq\operatorname{Ext}_{\gamma_{0,i}}\bigr),
 \label{eq:pdfv2-terminal-semilla-e}\\
 \operatorname{Upd}^{\mathsf E}_{\varepsilon}\mathsf E_{k,j}
 &:=
 \mathsf E_{k,j}\star
 \bigl(\varepsilon,\operatorname{Ext}_{\varepsilon},
 \operatorname{Ext}_{\varepsilon}\circ
 \operatorname{Ext}_{\gamma_{0,j}}
 \subseteq\operatorname{Ext}_{\gamma_{0,j+1}}\bigr).
 \label{eq:pdfv2-terminal-update-e}
\end{align}
L'identité et l'émission neutre prouvent la non-vacuité ; l'associativité de la composition relationnelle prouve la naturalité.

\paragraph{Relations \(\mathsf R\) : incidence, orientation et compatibilité de composition}
Soit \(\operatorname{RelBlock}_{j,r}\) le type des blocs d'équations étiquetés par l'arête et le préfixe sur lesquels ils sont vérifiés, et soit \(\operatorname{Pres}_{j,r}:=\operatorname{List}
(\operatorname{RelBlock}_{j,r})\). Le bloc unité et la graine sont
\begin{equation}
\label{eq:pdfv2-terminal-semilla-r}
 \mathsf B_I:=(I_{x_0};U_I=I,\ \mathsf U_I=I,\ \kappa(I)=0,\ \partial I=0),
 \qquad
 \mathsf R_{k,k}:=(\mathsf B_I)\star
 \mathop{\star}_{i=1}^{k}
 \mathsf B_R(\varepsilon_i;\gamma_{0,i-1}).
\end{equation}
Pour une arête \(\varepsilon:x_j\to x_{j+1}\), le nouveau bloc est le tuple étiqueté
\begin{equation}
\label{eq:pdfv2-terminal-bloque-r}
 \begin{split}
 \mathsf B_R(\varepsilon;\gamma_{0,j}):=\bigl(
 &\varepsilon,\gamma_{0,j};\ 
 \delta^\diamond(\varepsilon)
 =\operatorname{res}^\diamond(\varepsilon)
  +9\operatorname{quo}^\diamond(\varepsilon),\\
 &\Delta(\varepsilon)=\operatorname{or}(\varepsilon)
  +3\kappa(\varepsilon),\quad
 U_{\gamma_{0,j+1}}=U_\varepsilon U_{\gamma_{0,j}},\\
 &\mathsf U_{\gamma_{0,j+1}}
 =\mathsf U_\varepsilon\mathsf U_{\gamma_{0,j}},\quad
 \kappa(\gamma_{0,j+1})
 =\kappa(\varepsilon)+\mathsf C_\varepsilon\kappa(\gamma_{0,j}),\\
 &\partial\gamma_{0,j+1}
 =\partial\varepsilon+\partial\gamma_{0,j},\quad
 \operatorname{Cpl}_\varepsilon
 \text{ de }\eqref{eq:pdfv2-cor-acoplamientos-unicos}
 \bigr).
 \end{split}
\end{equation}
L'actualisation concatène le bloc sans perdre ni date ni multiplicité :
\begin{equation}
\label{eq:pdfv2-terminal-update-r}
 \operatorname{Upd}^{\mathsf R}_{\varepsilon}\mathsf R_{k,j}
 :=\mathsf R_{k,j}\star\mathsf B_R(\varepsilon;\gamma_{0,j}).
\end{equation}
La troncature \(\operatorname{tr}^{\mathsf R}_{j+1,j}\) supprime le dernier bloc étiqueté. En outre, \(a_*\mathsf B_R(\varepsilon;\gamma)
=\mathsf B_R(a\varepsilon;a\gamma)\), car chaque équation est naturelle sur la même arête et le même préfixe.

\paragraph{Nombres \(\mathsf N\) : expression de la valeur avec conservation de la généalogie}
Pour une arête \(\varepsilon:x_j\to x_{j+1}\), nous définissons la donnée numérique
\begin{equation}
\label{eq:pdfv2-terminal-numero-arista}
 \nu(\varepsilon):=
 \bigl(j+1,\mathsf Q_{j+1},g_{j+1},q^{(9)}_{j+1},\beta_{j+1},
 \operatorname{res}^{\Sigma},\operatorname{quo}^{\Sigma},
 \operatorname{res}^{\Pi},\operatorname{quo}^{\Pi},
 \operatorname{or},\kappa,\Delta\bigr)(\varepsilon).
\end{equation}
Son type est
\[
 \mathbb N\times\widetilde{\mathcal O}_9\times C_9\times\mathbb N
 \times\mathbb Z/12\mathbb Z
 \times\{0,\ldots,8\}^2\times\mathbb Z^5.
\]
La graine contient la liste numérique complète du préfixe \(0\to k\) ; l'actualisation ajoute exactement la donnée de l'arête suivante :
\begin{equation}
\label{eq:pdfv2-terminal-update-n}
 \mathsf N_{k,k}:=
 ((0,\mathsf Q_0,g_0,q^{(9)}_0,\beta_0;0,\ldots,0))
 \star\mathop{\star}_{i=1}^{k}\nu(\varepsilon_i),\qquad
 \operatorname{Upd}^{\mathsf N}_{\varepsilon}\mathsf N_{k,j}
 :=\mathsf N_{k,j}\star\nu(\varepsilon).
\end{equation}
La totalité provient des divisions APP, de la division TRIT et de l'odomètre TPK.

\paragraph{Orientation \(\mathsf O\) : transport du signe et sélection du feuillet}
Le type d'un état d'orientation est
\[
 \mathbb Z\times
 (\mathsf H^\Sigma_x\oplus\mathsf H^\Pi_x)
 \times\mathbb Z[\mathsf A],
\]
avec pour graine et actualisation
\begin{align}
 \mathsf O_{k,k}&:=
 \left(\sum_{i=1}^{k}\operatorname{or}(\varepsilon_i),
       \operatorname{Sh}(x_k),\partial\gamma_{0,k}\right),
 \label{eq:pdfv2-terminal-semilla-o}\\
 \operatorname{Upd}^{\mathsf O}_{\varepsilon}
 (\mathfrak o,\operatorname{Sh},\mathfrak b)
 &:=
 \bigl(\mathfrak o+\operatorname{or}(\varepsilon),
 U_\varepsilon\operatorname{Sh},
 \mathfrak b+\partial\varepsilon\bigr).
 \label{eq:pdfv2-terminal-update-o}
\end{align}
Les feuillets APP prouvent la non-vacuité. La fonctorialité de \(U\), l'additivité de \(\operatorname{or}\) et la télescopie de la frontière prouvent la naturalité.

\paragraph{Retenue \(\mathsf K\) : actualisation du quotient et mémoire résiduelle}
Nous utilisons les relèvements entiers typés dans \eqref{eq:pdfv2-cor-carry-matricial-def}. Le champ est de type
\[
 \mathsf M_{x_j}\times
 \operatorname{Hom}(\mathsf M_{x_0},\mathsf M_{x_j})
 \times\operatorname{Mat}_5(\mathbb Z)
 \times\operatorname{List}(\operatorname{Mat}_5(\mathbb Z))
 \times\operatorname{List}(\mathbb Q_{\geq0}).
\]
Pour \(\gamma_{k,j}=\varepsilon_j\cdots\varepsilon_{k+1}\), nous fixons explicitement
\begin{equation}
\label{eq:pdfv2-terminal-lift-entero-ruta}
 \widetilde L_{\gamma_{k,k}}:=I_5,\qquad
 \widetilde L_{\gamma_{k,j}}
 :=\widetilde L_r(\varepsilon_j)\cdots
   \widetilde L_r(\varepsilon_{k+1}).
\end{equation}
Pour le préfixe hérité, nous fixons également
\[
 \widetilde L_{\gamma_{0,k}}
 :=\widetilde L_r(\varepsilon_k)\cdots
   \widetilde L_r(\varepsilon_1).
\]
La graine et l'actualisation sont
\begin{align}
 \mathsf K_{k,k}&:=
 \bigl(\kappa(\gamma_{0,k}),\mathsf C_{\gamma_{0,k}},
       \widetilde L_{\gamma_{0,k}},
       \mathbf c_{\leq k},\mathbf k_{\leq k}\bigr),
 \label{eq:pdfv2-terminal-semilla-k}\\
 \operatorname{Upd}^{\mathsf K}_{\varepsilon}
 (\kappa_\gamma,\mathsf C_\gamma,\widetilde L_\gamma,
   \mathbf c,\mathbf k)
 &:=
 \bigl(\kappa(\varepsilon)+\mathsf C_\varepsilon\kappa_\gamma,\,
 \mathsf C_\varepsilon\mathsf C_\gamma,\,
 \widetilde L_r(\varepsilon)\widetilde L_\gamma,\,
 \mathbf c\star c_r(\varepsilon,\gamma),\,
 \mathbf k\star\kappa_j^{\rm can}\bigr).
 \label{eq:pdfv2-terminal-update-k}
\end{align}
L'identité fournit la graine ; les deux cocycles \eqref{eq:continuo-comun-cociclo-carry} et \eqref{eq:pdfv2-cor-cociclo-matricial} prouvent l'indépendance du parenthésage et la naturalité.

\paragraph{Mémoire \(\mathsf M\) : registre du parcours et compatibilité sous troncature}
Son type conserve simultanément mémoire affine, registre et mémoires signées :
\begin{equation}
\label{eq:pdfv2-terminal-tipo-ledger-entry}
 \operatorname{LedgerEntry}:=
 \{\lambda(\varepsilon):\varepsilon\text{ émission élémentaire typée}\}.
\end{equation}
\[
 \mathsf M_x\times\operatorname{List}(\operatorname{LedgerEntry})
 \times\mathbb N^2\times\mathbb Z\times\mathbb N.
\]
Nous posons
\begin{align}
 \mathsf M_{k,k}&:=
 \bigl(m_k,\operatorname{Led}(\gamma_{0,k}),
       M_k^+,M_k^-,D_k,H_k\bigr),
 \label{eq:pdfv2-terminal-semilla-m}\\
 \operatorname{Upd}^{\mathsf M}_{\varepsilon}
 (m,\mathbf{Led},M^+,M^-,D,H)
 &:=
 \bigl(\mathsf C_\varepsilon m+\kappa(\varepsilon),\,
 \mathbf{Led}\star\lambda(\varepsilon),\,
 M^++b_\varepsilon^+,\,M^-+b_\varepsilon^-,\,
 D+b_\varepsilon,\,H+|b_\varepsilon|\bigr).
 \label{eq:pdfv2-terminal-update-m}
\end{align}
L'action affine et \eqref{eq:pdfv2-cor-memorias} prouvent la totalité et la naturalité ; le passé n'est pas recalculé à partir de la dernière valeur.

\paragraph{Frontière \(\mathsf B\) : données de raccordement et survie de l'extension}
Le type est une chaîne de frontière accompagnée d'un mot de complexes munis de bases et d'applications de chaînes. Sa graine et son actualisation sont
\[
 \mathsf B_{k,j}\in
 \operatorname{Rec}[\mathfrak b\in\mathbb Z[\mathsf A_j];
 C_{x_j,r}\in\operatorname{Ch}_{\rm bas};
 \mathbf F\in\operatorname{Path}_{\operatorname{Ch}_{\rm bas}}
 (C_{x_j,r},C_{x_0,r})].
\]
\begin{align}
 \mathsf B_{k,k}&:=
 \bigl(\partial\gamma_{0,k},C_{x_k,r},\mathbf F_{\leq k}\bigr),
 \label{eq:pdfv2-terminal-semilla-b}\\
 \operatorname{Upd}^{\mathsf B}_{\varepsilon}
 (\mathfrak b,C_{x_j,r},\mathbf F)
 &:=
 \bigl(\mathfrak b+\partial\varepsilon,\,
 C_{x_{j+1},r},\,\mathbf F\star F_\varepsilon\bigr).
 \label{eq:pdfv2-terminal-update-b}
\end{align}
La non-vacuité provient des ports communs ; \(\partial^2=0\), \(F_{\beta\varepsilon}=F_\varepsilon F_\beta\) et \eqref{eq:continuo-comun-frontera-telescopica} prouvent la compatibilité.

\paragraph{Parcours \(\mathsf G\) : histoire observable de l'extension survivante}
Le champ est la liste des préfixes dans la catégorie des chemins :
\begin{equation}
\label{eq:pdfv2-terminal-update-g}
 \mathsf G_{k,k}:=(\gamma_{0,i})_{0\leq i\leq k},\qquad
 \operatorname{Upd}^{\mathsf G}_{\varepsilon}\mathsf G_{k,j}
 :=\mathsf G_{k,j}\star\gamma_{0,j+1}.
\end{equation}
L'identité, la concaténation associative et \(a(\varepsilon\gamma)=a(\varepsilon)a(\gamma)\) prouvent les trois propriétés requises.

\paragraph{Multisection \(\mathsf X\) : fibre finie stable et obstruction à la sélection ponctuelle}
Pour \(M\geq j\), nous définissons par action
\begin{equation}
\label{eq:pdfv2-terminal-multiseccion-interna}
 \mathsf X_{k,j}(M):=
 \{\gamma_{k,j}\star q:
 q\in\operatorname{Term}_{j,M}(x_j)\}.
\end{equation}
Ainsi
\[
 \mathsf X_{k,k}(M)=\operatorname{Term}_{k,M}(x_k),
\]
et
\begin{equation}
\label{eq:pdfv2-terminal-update-x}
 \operatorname{Upd}^{\mathsf X}_{\varepsilon}\mathsf X_{k,j}(M)
 :=
 \{\gamma_{k,j}\star\varepsilon\star q:
 q\in\operatorname{Term}_{j+1,M}(x_{j+1})\}.
\end{equation}
Le prolongement neutre prouve la non-vacuité. L'effacement du dernier pas et l'action d'une symétrie commutent littéralement avec la concaténation.

\paragraph{Survie \(\mathsf S\) : persistance sous raffinement et passage à la limite}
Ici, on n'enregistre pas une proposition, mais des témoins. Soit
\begin{equation}
\label{eq:pdfv2-terminal-testigo-neutro}
 s_{j,j}:=I_{x_j},\qquad
 s_{j,M+1}:=s_{j,M}\star
 \varepsilon_M^0(t(s_{j,M})).
\end{equation}
La graine est \(\mathsf S_{k,k}:=(s_{k,M})_{M\geq k}\), et
\begin{equation}
\label{eq:pdfv2-terminal-update-s}
 \operatorname{Upd}^{\mathsf S}_{\varepsilon}\mathsf S_{k,j}
 :=(\gamma_{k,j}\star\varepsilon\star s_{j+1,M})_{M\geq j+1}.
\end{equation}
Par construction, \(\pi_{M+1,M}s_{j,M+1}=s_{j,M}\) ; la naturalité de l'émission neutre prouve l'équivariance.

\paragraph{Terminal \(\mathsf T\) : universalité de l'objet limite et factorisation des lecteurs}
Le terminal est le système inverse, non une queue choisie :
\begin{equation}
\label{eq:pdfv2-terminal-sistema-inverso-campo}
 \mathsf T(x_j):=
 \bigl((\operatorname{Term}_{j,M}(x_j))_{M\geq j},
       (\pi_{M+1,M})_{M\geq j}\bigr).
\end{equation}
Pour chaque \(M\geq j\), l'application de préfixe typée est
\begin{equation}
\label{eq:pdfv2-terminal-update-t}
 \operatorname{pref}_{\gamma_{k,j},M}:
 \operatorname{Term}_{j,M}(x_j)\longrightarrow
 \operatorname{Term}_{k,M}(x_k),\qquad
 q\longmapsto\gamma_{k,j}\star q.
\end{equation}
Nous définissons l'état du champ et son actualisation par
\begin{align}
 \mathsf T_{k,j}
 &:=
 \left(\mathsf T(x_j),
   (\operatorname{pref}_{\gamma_{k,j},M})_{M\geq j}\right),
 \label{eq:pdfv2-terminal-estado-t}\\
 \operatorname{Upd}^{\mathsf T}_{\varepsilon}\mathsf T_{k,j}
 &:=
 \left(\mathsf T(x_{j+1}),
   (\operatorname{pref}_{\gamma_{k,j+1},M})_{M\geq j+1}\right),
 \qquad \gamma_{k,j+1}:=\gamma_{k,j}\star\varepsilon.
 \label{eq:pdfv2-terminal-update-t-explicita}
\end{align}
En particulier, \(\mathsf T_{k,k}\) utilise \(\gamma_{k,k}=I_{x_k}\). Les témoins \eqref{eq:pdfv2-terminal-testigo-neutro} et la surjectivité de \(\pi\) prouvent que le système est non vide ; niveau par niveau,
\begin{equation}
\label{eq:pdfv2-terminal-prefijo-natural}
 \pi^{\mathscr H}_{M+1,M}
 \operatorname{pref}_{\gamma_{k,j},M+1}
 =\operatorname{pref}_{\gamma_{k,j},M}
  \pi^{\mathscr H}_{M+1,M},
\end{equation}
car les deux parcours suppriment seulement la dernière émission de \(q\).

\paragraph{13. Lecteur terminal \(\mathsf L\) : évaluation de l'extension compatible et conservation de la généalogie}
Le lecteur ne contient aucun objet classique. Son type est la catégorie des chemins de l'unique constructeur dépendant \(\mathfrak D_{j,r}\) de \eqref{eq:pdfv2-cor-constructor-unico}. Nous posons
\begin{equation}
\label{eq:pdfv2-terminal-transicion-lector-tipado}
 u^{\mathfrak D}_{j+1,j}
 :=u_{(j+1,r),(j,r)}:
 \mathfrak D_{j+1,r}(x_{j+1})
 \longrightarrow\mathfrak D_{j,r}(x_j),
\end{equation}
la composante du transport commun \eqref{eq:pdfv2-cor-naturalidad-unica} qui tronque la dernière émission, réduit les coefficients et précompose les ports. Par composition,
\begin{equation}
\label{eq:pdfv2-terminal-transicion-lector-composicion}
 u^{\mathfrak D}_{m,k}
 =u^{\mathfrak D}_{\ell,k}u^{\mathfrak D}_{m,\ell}
 \qquad(m\geq\ell\geq k).
\end{equation}
\begin{align}
 \mathsf L_{k,k}
 &:=(\mathfrak D_{0,r}(x_0))\star
 \mathop{\star}_{i=1}^{k}
 \bigl(\mathfrak D_{i,r}(x_i),u^{\mathfrak D}_{i,i-1},
       \mathcal O(\mathfrak e_{\varepsilon_i})\bigr),
 \label{eq:pdfv2-terminal-semilla-l}\\
 \operatorname{Upd}^{\mathsf L}_{\varepsilon}\mathsf L_{k,j}
 &:=
 \mathsf L_{k,j}\star
 \bigl(\mathfrak D_{j+1,r}(x_{j+1}),
 u^{\mathfrak D}_{j+1,j},\mathcal O(\mathfrak e_\varepsilon)\bigr).
 \label{eq:pdfv2-terminal-update-l}
\end{align}
La totalité de \(\mathfrak D\) et \eqref{eq:pdfv2-terminal-transicion-lector-composicion} prouvent la non-vacuité et la naturalité.

Les treize graines précédentes forment un élément unique
\begin{equation}
\label{eq:pdfv2-terminal-semilla-trece}
 \mathfrak r^0_{k,r}:=
 \langle
 \mathsf F_{k,k};\mathsf E_{k,k};\mathsf R_{k,k};\mathsf N_{k,k};
 \mathsf O_{k,k};\mathsf K_{k,k};\mathsf M_{k,k};\mathsf B_{k,k};
 \mathsf G_{k,k};\mathsf X_{k,k};\mathsf S_{k,k};\mathsf T_{k,k};
 \mathsf L_{k,k}
 \rangle\in\mathfrak R_{k,r}.
\end{equation}
L'action d'une arête est l'application totale
\begin{equation}
\label{eq:pdfv2-terminal-actualizacion-unica}
 \operatorname{Upd}_\varepsilon:
 \mathfrak R_{j,r}(\gamma_{0,j})
 \longrightarrow\mathfrak R_{j+1,r}(\gamma_{0,j+1}),\qquad
 \operatorname{Upd}_\varepsilon:=
 \langle\operatorname{Upd}^{\mathsf F}_\varepsilon,\ldots,
 \operatorname{Upd}^{\mathsf L}_\varepsilon\rangle,
\end{equation}
avec les treize actions \eqref{eq:pdfv2-terminal-update-f}--\eqref{eq:pdfv2-terminal-update-l}. Enfin, le registre terminal est \emph{engendré} par le pliage
\begin{equation}
\label{eq:pdfv2-terminal-registro-trece}
 \boxed{
 \mathfrak r_{k,N}(h):=
 \operatorname{Upd}_{\varepsilon_N}\circ\cdots\circ
 \operatorname{Upd}_{\varepsilon_{k+1}}(\mathfrak r^0_{k,r}).}
\end{equation}
Ce n'est pas le dépaquetage d'une copie intacte de l'état. Sur l'image de ce pliage, on définit, composante par composante, la troncature
\begin{equation}
\label{eq:pdfv2-terminal-truncamiento-trece-def}
 \operatorname{tr}_{13;j+1,j}
 \mathfrak r_{k,j+1}(h\star\varepsilon)
 :=\mathfrak r_{k,j}(h).
\end{equation}
Elle est bien définie : les champs \(\mathsf E\) et \(\mathsf G\) conservent la liste ordonnée des émissions et des préfixes, si bien que l'égalité de deux registres implique l'égalité de leurs histoires étiquetées. Dans \(\mathsf F,\mathsf E,\mathsf R,\mathsf N,\mathsf G,\mathsf T,\mathsf L\), la dernière entrée ou le dernier bloc est supprimé ; dans \(\mathsf O,\mathsf K,\mathsf M,
\mathsf B\), la formule cumulative est appliquée au préfixe retrouvé ; et dans \(\mathsf X,\mathsf S\), le cône des queues est restreint. On n'inverse pas une union idempotente et l'on ne sélectionne pas une continuation.

\begin{proposition}[Totalité et naturalité des treize actualisations]
\label{prop:pdfv2-terminal-trece-total-natural}
Toutes les graines de \eqref{eq:pdfv2-terminal-semilla-trece} existent, chaque \(\operatorname{Upd}_\varepsilon\) est total et, pour toute symétrie cohérente \(a\) et toute troncature d'horizon,
\begin{equation}
\label{eq:pdfv2-terminal-trece-naturalidad}
 a_*\operatorname{Upd}_\varepsilon
 =\operatorname{Upd}_{a\varepsilon}a_*,
 \qquad
 \operatorname{tr}_{13;j+1,j}
 \operatorname{Upd}_\varepsilon\mathfrak r_{k,j}(h)
 =\mathfrak r_{k,j}(h).
\end{equation}
\end{proposition}

\begin{proof}
L'existence de chaque graine a été exposée dans les treize paragraphes. Dans \(\mathsf F,\mathsf E,\mathsf G,\mathsf X,\mathsf S,\mathsf T,\mathsf L\), la formule est une concaténation avec une application déjà typée. Dans \(\mathsf R,\mathsf N,\mathsf O,\mathsf K,\mathsf M,\mathsf B\), la formule est, respectivement, la concaténation d'un bloc d'équations étiqueté, l'adjonction d'une donnée numérique, le transport de feuillets, la composition de cocycles, l'action affine et l'application de chaînes. Les identités de naturalité citées après chaque formule prouvent la première égalité de \eqref{eq:pdfv2-terminal-trece-naturalidad} composante par composante. La seconde est exactement \eqref{eq:pdfv2-terminal-truncamiento-trece-def} ; sa bonne définition utilise l'histoire ordonnée conservée par \(\mathsf E\) et \(\mathsf G\).
\end{proof}

\begin{proposition}[Certificat conjoint des cinq développements internes]
\label{prop:pdfv2-terminal-cinco-desarrollos-integrales}
Les actualisations des treize coordonnées transportent conjointement : la télescopie spectrale--polaire et le retour sans réinitialisation de \cref{thm:continuo-sp-telescopia-no-autonoma,eq:continuo-sp-no-reset} ; le bilan orthogonal, la colonne d'observation et l'identité de Stein de \cref{eq:pdfv2-sol-ortogonalidad,eq:pdfv2-sol-stein} ; la naturalité directe et adjointe ainsi que le rétracte cellulaire de \cref{eq:gauint-naturalidad-directa-adjunta,eq:gauint-sdr} ; la totalisation, le cône du retour nonadique et les fibres complètes de \cref{eq:pdfv2-coh-total-square-zero,eq:pdfv2-coh-terminal-cone,eq:pdfv2-coh-extension-multisection} ; et la diagonale d'Alexander--Whitney, la réduction de Smith et la tour de Bockstein de \cref{eq:detint-aw-natural,eq:detint-smith,eq:detint-bock-leading}. Toutes ces actions ont pour domaine le même registre engendré par \eqref{eq:pdfv2-terminal-registro-trece} ; aucune n'introduit de sous-domaine initial, d'histoire privilégiée ou de sélection motivée par une application ultérieure.
\end{proposition}

\begin{proof}
Chaque identité citée agit sur une coordonnée du registre commun, et ses applications de troncature sont des restrictions de \eqref{eq:pdfv2-terminal-truncamiento-trece-def}. La naturalité composante par composante démontrée dans \eqref{eq:pdfv2-terminal-trece-naturalidad} permet de les assembler en une seule actualisation. L'égalité des domaines et la conservation de l'histoire ordonnée empêchent de remplacer la génération corrélative par cinq projections rétrospectives.
\end{proof}

\section{Terminal multivalué sans sélecteur}

Le terminal fini engendré par \(\widehat x_k\) est la multisection
\begin{equation}
\label{eq:pdfv2-terminal-multivaluado}
 \mathfrak T_{k,N}(\widehat x_k)
 :=\left\{\mathfrak r_{k,N}(h):
            h\in\mathscr H_{k,N}(\widehat x_k)\right\}.
\end{equation}
L'expression du membre de droite renvoie l'ensemble entier. Elle ne contient aucune règle distinguant une histoire privilégiée.

L'espace ambiant terminal du niveau s'obtient en réunissant ces multisections sans identifier les registres :
\begin{equation}
\label{eq:pdfv2-terminal-ambiente-nivel}
 \mathfrak T_{k,N}
 :=\bigcup_{\widehat x_k\in\widehat{\mathcal X}^{\rm tot}_k}
   \mathfrak T_{k,N}(\widehat x_k).
\end{equation}

La version linéaire, lorsqu'il faut sommer des contributions, se définit dans le module libre par
\begin{equation}
\label{eq:pdfv2-terminal-evaluacion-todas}
 \begin{aligned}
 \mathsf E_{k,N}^{\rm all}:
 \mathbb Z[\widehat{\mathcal X}^{\rm tot}_k]&\longrightarrow
 \mathbb Z[\mathfrak T_{k,N}],\\
 [\widehat x_k]&\longmapsto
 \sum_{h\in\mathscr H_{k,N}(\widehat x_k)}
 [\mathfrak r_{k,N}(h)].
 \end{aligned}
\end{equation}
Chaque témoin apparaît avec sa multiplicité. S'il existe une mesure projective positive \(\mu\), la version isométrique utilise toutes les probabilités conditionnelles :
\begin{equation}
\label{eq:pdfv2-terminal-evaluacion-ponderada}
 \widehat{\mathsf E}_{k,N}^{\mu}e_{\widehat x_k}
 :=\sum_{h\in\mathscr H_{k,N}(\widehat x_k)}
 \mu_{\widehat x_k,N}(h)^{1/2}
 e_{\mathfrak r_{k,N}(h)}.
\end{equation}
La projectivité donne
\begin{equation}
\label{eq:pdfv2-terminal-isometria}
 \left\|\widehat{\mathsf E}_{k,N}^{\mu}e_{\widehat x_k}\right\|^2
 =\sum_{h\in\mathscr H_{k,N}(\widehat x_k)}
   \mu_{\widehat x_k,N}(h)=1.
\end{equation}
La somme n'est pas remplacée par ``l'histoire qui engendre le cylindre''.

\section{La connexion nonadique comme correspondance}

L'holonomie nonadique a été engendrée par l'odomètre et la connexion TPK, sans hypothèses d'appartenance, dans \eqref{eq:continuo-comun-gamma9-apice}--\eqref{eq:continuo-comun-gamma9-relacion}. Dans ce chapitre, nous utilisons
\begin{equation}
\label{eq:pdfv2-terminal-apice-gamma9}
 \mathsf Z_{9,k}:=\mathsf Z^{\Gamma}_{9,k}.
\end{equation}
Les applications
\begin{equation}
\label{eq:pdfv2-terminal-span-gamma9}
 \widehat{\mathcal X}^{\rm tot}_k
 \xleftarrow{\ s_{9,k}\ }
 \mathsf Z^{\Gamma}_{9,k}
 \xrightarrow{\ t_{9,k}\ }
 \widehat{\mathcal X}^{\rm tot}_{k+9}
\end{equation}
sont ceux de \eqref{eq:continuo-comun-gamma9-span} ; son application source est surjective par \cref{prop:continuo-comun-gamma9-total-natural}. Chaque point du sommet conserve les neuf émissions compatibles avec la phase. Il n'est pas remplacé par la paire de ses extrémités.

Pour \(N\geq k+9\), nous définissons le sommet d'horizon avant linéarisation :
\begin{equation}
\label{eq:pdfv2-terminal-apice-horizonte-gamma9}
 \mathsf Z^\Gamma_{9;k,N}:=
 \coprod_{z=(\widehat x_k,\boldsymbol\omega)
            \in\mathsf Z^\Gamma_{9,k}}
 \operatorname{Term}_{k+9,N}(t_{9,k}z).
\end{equation}
Un élément s'écrit \((z,q)\), où \(\boldsymbol\omega\) est le mot nonadique de neuf émissions et \(q\) sa queue jusqu'à \(N\). Les applications sur les histoires utilisent le codomaine étiqueté
\begin{equation}
\label{eq:pdfv2-terminal-gamma9-codominio-colas}
 \mathscr D^\Gamma_{9;k,N}:=
 \coprod_{z\in\mathsf Z^\Gamma_{9,k}}
 \operatorname{Term}_{k+9,N}(t_{9,k}z).
\end{equation}
Ainsi, les applications sont
\begin{align}
 S_{k,N}(z,q)&:=\boldsymbol\omega\star q,
 \label{eq:pdfv2-terminal-gamma9-s-horizonte}\\
 T_{k,N}(z,q)&:=(z,q)\in\mathscr D^\Gamma_{9;k,N}.
 \label{eq:pdfv2-terminal-gamma9-t-horizonte}
\end{align}
La première aboutit dans \(\coprod_x\mathscr H_{k,N}(x)\) ; la seconde conserve expressément le facteur étiqueté par \(z\).

L'application de liaison du sommet est l'action explicite
\begin{equation}
\label{eq:pdfv2-terminal-gamma9-enlace-apice}
 \pi^Z_{N+1,N}(z,q\star\varepsilon):=(z,q).
\end{equation}
Elle est surjective : étant donné \((z,q)\), on ajoute à la queue l'émission neutre de l'extrémité de \(q\). Le codomaine des queues possède la liaison
\begin{equation}
\label{eq:pdfv2-terminal-gamma9-enlace-codominio}
 \pi^{\mathscr D}_{N+1,N}(z,q\star\varepsilon):=(z,q).
\end{equation}
Nous notons \(\pi^{\mathscr H,j}_{N+1,N}\) l'application \eqref{eq:pdfv2-terminal-truncamiento-historias} de niveau initial \(j\). Alors
\begin{align}
 \pi^{\mathscr H,k}_{N+1,N}S_{k,N+1}
 &=S_{k,N}\pi^Z_{N+1,N},
 \label{eq:pdfv2-terminal-gamma9-cuadrado-fuente}\\
 \pi^{\mathscr D}_{N+1,N}T_{k,N+1}
 &=T_{k,N}\pi^Z_{N+1,N}.
 \label{eq:pdfv2-terminal-gamma9-cuadrado-destino}
\end{align}
Dans le premier carré, les deux parcours envoient \(\boldsymbol\omega\star q\star\varepsilon\) sur \(\boldsymbol\omega\star q\) ; dans le second, ils envoient \((z,q\star\varepsilon)\) sur \((z,q)\). Par décomposition unique d'un mot en ses neuf premières émissions et sa queue, \(S_{k,N}\) est bijectif. Le second carré est cartésien par la définition en coproduit fibré de \eqref{eq:pdfv2-terminal-apice-horizonte-gamma9}. Ce sont les carrés de Beck--Chevalley qui conservent témoins et multiplicités.

Le sommet correspondant des registres, désormais engendré par les treize actualisations, est
\begin{equation}
\label{eq:pdfv2-terminal-apice-terminal-gamma9}
 \begin{split}
 \mathsf Z^{\mathfrak T}_{9;k,N}:=\bigl\{
 (&z,q;
   \mathfrak r_{k,N}(\boldsymbol\omega\star q),
   \mathfrak r_{k+9,N}(q)):\\
 &z=(\widehat x_k,\boldsymbol\omega)\in
       \mathsf Z^\Gamma_{9,k},\quad
 q\in\operatorname{Term}_{k+9,N}(t_{9,k}z)
 \bigr\}.
 \end{split}
\end{equation}
Ses applications source et cible lisent respectivement le premier ou le second registre, mais ne suppriment jamais de leur domaine le témoin \((z,q)\) :
\begin{equation}
\label{eq:pdfv2-terminal-span-terminal}
 \mathfrak T_{k,N}
 \xleftarrow{\ s^{\mathfrak T}_{9;k,N}\ }
 \mathsf Z^{\mathfrak T}_{9;k,N}
 \xrightarrow{\ t^{\mathfrak T}_{9;k,N}\ }
 \mathfrak T_{k+9,N}.
\end{equation}
Si \(\zeta=(z,q;r_s,r_t)\), son action est
\begin{equation}
\label{eq:pdfv2-terminal-span-terminal-accion}
 s^{\mathfrak T}_{9;k,N}(\zeta):=r_s,
 \qquad
 t^{\mathfrak T}_{9;k,N}(\zeta):=r_t.
\end{equation}
Sa liaison est définie par
\begin{equation}
\label{eq:pdfv2-terminal-gamma9-enlace-registros}
 \begin{split}
 \pi^{Z,\mathfrak T}_{N+1,N}
 \bigl(&z,q\star\varepsilon;
       \mathfrak r_{k,N+1}
          (\boldsymbol\omega\star q\star\varepsilon),
       \mathfrak r_{k+9,N+1}(q\star\varepsilon)\bigr)\\
 &:=(z,q;\mathfrak r_{k,N}(\boldsymbol\omega\star q),
       \mathfrak r_{k+9,N}(q)).
 \end{split}
\end{equation}
La paire \((z,q)\) est une coordonnée du sommet, non une représentation choisie d'un triplet ; la liaison est donc bien définie même si deux registres terminaux coïncident. La compatibilité de \(\operatorname{Upd}\) avec l'effacement de la dernière arête prouve que \eqref{eq:pdfv2-terminal-gamma9-enlace-registros} commute avec les deux applications de \eqref{eq:pdfv2-terminal-span-terminal}.

Si l'on souhaite une action sur des modules libres, cette correspondance est linéarisée en sommant tous les témoins :
\begin{equation}
\label{eq:pdfv2-terminal-linealizacion-span}
 [\Gamma_9]_{k,N}[r]
 :=\sum_{\zeta\in(s^{\mathfrak T}_{9;k,N})^{-1}(r)}
 [t^{\mathfrak T}_{9;k,N}(\zeta)].
\end{equation}
La formule conserve les multiplicités. L'égalité des phases après neuf pas ne contracte ni les états, ni les témoins, ni leurs registres.

\section{Unicité de la décomposition par préfixes et naturalité du système de troncatures}

Pour \(k\leq\ell\leq N\), toute histoire de \(\mathscr H_{k,N}(\widehat x_k)\) possède un unique préfixe \(p\in\mathscr P_{\ell\mid k}(\widehat x_k)\). Il existe donc une décomposition disjointe par préfixes, exprimée sans en choisir aucun :
\begin{equation}
\label{eq:pdfv2-terminal-particion-historias}
 \mathscr H_{k,N}(\widehat x_k)
 =\bigsqcup_{p\in\mathscr P_{\ell\mid k}(\widehat x_k)}
 \left\{p\star q:q\in\mathscr H_{\ell,N}(p)\right\}.
\end{equation}
Soit \(\operatorname{Conc}_p\) l'opération qui place en tête le registre du préfixe \(p\) et applique successivement \eqref{eq:pdfv2-terminal-actualizacion-unica}. Nous définissons la transition commune sur un système complet de terminaux par
\begin{equation}
\label{eq:pdfv2-terminal-transicion-comun}
 \mathfrak u_{\ell,k}
 \left(
   \bigl(\mathfrak T_{\ell,N}(p)\bigr)_{p\in\mathscr P_{\ell\mid k}(\widehat x_k)}
 \right)
 :=\bigcup_{p\in\mathscr P_{\ell\mid k}(\widehat x_k)}
   \operatorname{Conc}_p\mathfrak T_{\ell,N}(p).
\end{equation}

En écrivant \(\mathbf G_{k,N}(\widehat x_k):=\mathfrak T_{k,N}(\widehat x_k)\), toute la naturalité se concentre dans l'identité unique
\begin{equation}
\label{eq:pdfv2-terminal-naturalidad-unica}
 \boxed{
 \mathfrak u_{\ell,k}
 \left(
   \bigl(\mathbf G_{\ell,N}(p)\bigr)_{p\in\mathscr P_{\ell\mid k}(\widehat x_k)}
 \right)
 =\mathbf G_{k,N}(\widehat x_k).}
\end{equation}
Il n'existe pas de loi distincte pour chaque discriminant interne. Les cinq caractéristiques obéissent à \eqref{eq:pdfv2-terminal-naturalidad-unica} puisque chacune s'actualise au moyen du même préfixe, de la même extension et du même registre.

La correspondance nonadique est compatible avec cette identité au moyen des applications déjà construites, non d'une condition implicite. À tout horizon, \eqref{eq:pdfv2-terminal-gamma9-cuadrado-fuente} et \eqref{eq:pdfv2-terminal-gamma9-cuadrado-destino} conservent le témoin nonadique avec sa multiplicité, et \eqref{eq:pdfv2-terminal-gamma9-enlace-registros} fait de même dans les treize coordonnées. Si \(a\) est une symétrie cohérente de l'état commun, son action
\[
 a^Z(z,q):=(a^Zz,aq)
\]
satisfait
\begin{equation}
\label{eq:pdfv2-terminal-gamma9-equivarianza-completa}
 S a^Z=aS,\qquad T a^Z=aT,\qquad
 \pi^Za^Z=a^Z\pi^Z,\qquad
 s^{\mathfrak T}a^Z=as^{\mathfrak T},\qquad
 t^{\mathfrak T}a^Z=at^{\mathfrak T}.
\end{equation}
Ces égalités se vérifient entrée par entrée au moyen de \eqref{eq:continuo-comun-gamma9-naturalidad} et de la naturalité des treize actualisations. Elles constituent la condition concrète permettant d'utiliser \eqref{eq:pdfv2-terminal-linealizacion-span} ; une simple égalité de valeurs visibles ne la remplace pas.

\section{Passage à la limite sans intervertir image et limite}

La troncature d'une histoire et de chacun des treize champs définit
\begin{equation}
\label{eq:pdfv2-terminal-transicion-horizonte}
 \pi^{\mathfrak T}_{N+1,N}:
 \mathfrak T_{k,N+1}(\widehat x_k)\longrightarrow
 \mathfrak T_{k,N}(\widehat x_k),
 \qquad
 \pi^{\mathfrak T}_{N+1,N}\mathfrak r_{k,N+1}(h)
 =\mathfrak r_{k,N}(h_{\leq N}).
\end{equation}
Elle est surjective par \eqref{eq:pdfv2-terminal-truncamiento-historias}. Le terminal complet est alors
\begin{equation}
\label{eq:pdfv2-terminal-limite-inverso}
 \mathfrak T_{k,\infty}(\widehat x_k)
 :=\left\{(r_N)_{N\geq k}:
 r_N\in\mathfrak T_{k,N}(\widehat x_k),\quad
 \pi^{\mathfrak T}_{N+1,N}r_{N+1}=r_N\right\}.
\end{equation}
Les ensembles finis de \eqref{eq:pdfv2-terminal-multivaluado} sont non vides et leurs transitions sont surjectives ; ainsi
\begin{equation}
\label{eq:pdfv2-terminal-limite-no-vacio}
 \mathfrak T_{k,\infty}(\widehat x_k)\neq\varnothing.
\end{equation}

\begin{corollary}[Persistance simultanée des cinq développements]
\label{cor:pdfv2-terminal-limite-cinco-desarrollos}
Dans chaque élément de \(\mathfrak T_{k,\infty}(\widehat x_k)\) persistent simultanément la limite forte \(\mathcal T_\infty\) de la télescopie polaire, le système de Gram et de pertes, la mesure d'incidence avec ses deux naturalités, le cône \(\operatorname{Cone}(I-U_{\Gamma_9,N})\) et le système compatible des formes de Smith et des pages de Bockstein. Aucune de ces données ne s'obtient en reconstruisant la seule valeur terminale visible.
\end{corollary}

\begin{proof}
La surjectivité des troncatures conserve les témoins complets. Les équations de télescopie, de naturalité et de réduction de \cref{prop:pdfv2-terminal-cinco-desarrollos-integrales} commutent avec ces troncatures ; par universalité de la limite inverse, elles définissent un unique système compatible sur chaque histoire survivante.
\end{proof}

Le sommet nonadique à la limite est défini, non inféré d'une image :
\begin{equation}
\label{eq:pdfv2-terminal-gamma9-apice-limite-raw}
 \mathsf Z^\Gamma_{9;k,\infty}
 :=\varprojlim_{N\geq k+9}
 \bigl(\mathsf Z^\Gamma_{9;k,N},\pi^Z_{N+1,N}\bigr).
\end{equation}
Puisque les liaisons ne modifient pas le témoin fini \(z\) et tronquent seulement sa queue, il existe l'isomorphisme explicite
\begin{equation}
\label{eq:pdfv2-terminal-gamma9-apice-limite-msect}
 \mathsf Z^\Gamma_{9;k,\infty}
 \xrightarrow{\ \cong\ }
 \coprod_{z\in\mathsf Z^\Gamma_{9,k}}
 \operatorname{Msect}(t_{9,k}z),\qquad
 ((z,q_N))_N\longmapsto(z,(q_N)_N).
\end{equation}
L'inverse prend un témoin \(z\) et une queue compatible \((q_N)_N\), puis forme \(((z,q_N))_N\). Les deux compositions sont des identités coordonnée par coordonnée. Les applications limites sont
\begin{equation}
\label{eq:pdfv2-terminal-gamma9-mapas-limite-raw}
 S_{k,\infty}((z_N)_N):=(S_{k,N}z_N)_N,\qquad
 T_{k,\infty}((z_N)_N):=(T_{k,N}z_N)_N;
\end{equation}
sont bien définies précisément par \eqref{eq:pdfv2-terminal-gamma9-cuadrado-fuente} et \eqref{eq:pdfv2-terminal-gamma9-cuadrado-destino}.

On définit de même le sommet des registres
\begin{equation}
\label{eq:pdfv2-terminal-gamma9-apice-registros-limite}
 \mathsf Z^{\mathfrak T}_{9;k,\infty}
 :=\varprojlim_{N\geq k+9}
 \bigl(\mathsf Z^{\mathfrak T}_{9;k,N},
       \pi^{Z,\mathfrak T}_{N+1,N}\bigr).
\end{equation}
Les correspondances \eqref{eq:pdfv2-terminal-span-terminal} forment alors, composante par composante, le span
\begin{equation}
\label{eq:pdfv2-terminal-span-limite}
 \mathfrak T_{k,\infty}
 \xleftarrow{\ s^{\mathfrak T}_{9;k,\infty}\ }
 \mathsf Z^{\mathfrak T}_{9;k,\infty}
 \xrightarrow{\ t^{\mathfrak T}_{9;k,\infty}\ }
 \mathfrak T_{k+9,\infty}.
\end{equation}
L'image d'une limite n'a pas été identifiée à la limite d'images : \eqref{eq:pdfv2-terminal-limite-inverso}, \eqref{eq:pdfv2-terminal-gamma9-apice-limite-raw} et \eqref{eq:pdfv2-terminal-gamma9-apice-registros-limite} construisent directement les trois systèmes compatibles et leurs applications.

\section{Reconstruction interne de l'état enrichi à partir de ses \(13\) coordonnées généalogiques}

La reconstruction ne dépaquette pas des noms : elle répète le pliage qui a engendré le registre. Pour \(h=(\varepsilon_{k+1},\ldots,\varepsilon_N)\), nous définissons
\begin{equation}
\label{eq:pdfv2-terminal-recuperacion-trece}
 \boxed{
 \operatorname{Rec}_{13;k,N}(h):=
 \operatorname{Upd}_{\varepsilon_N}\circ\cdots\circ
 \operatorname{Upd}_{\varepsilon_{k+1}}
 (\mathfrak r^0_{k,r}).}
\end{equation}
Par \eqref{eq:pdfv2-terminal-registro-trece}, \(\operatorname{Rec}_{13;k,N}(h)=\mathfrak r_{k,N}(h)\), mais l'égalité est désormais une identité entre deux compositions d'actions, non la projection d'une entrée préservée. Si \(\operatorname{pr}_{\mathsf A}\) désigne le projecteur du registre sur le champ \(\mathsf A\in\{\mathsf F,\mathsf E,\mathsf R,\mathsf N,\mathsf O,
\mathsf K,\mathsf M,\mathsf B,\mathsf G,\mathsf X,\mathsf S,\mathsf T,
\mathsf L\}\), alors
\begin{equation}
\label{eq:pdfv2-terminal-recuperacion-campo-fold}
 \operatorname{pr}_{\mathsf A}\operatorname{Rec}_{13;k,N}(h)
 =\operatorname{Upd}^{\mathsf A}_{\varepsilon_N}\circ\cdots\circ
  \operatorname{Upd}^{\mathsf A}_{\varepsilon_{k+1}}
  (\mathsf A_{k,k}).
\end{equation}
La preuve se fait par récurrence sur \(N-k\) : pour \(N=k\), les deux membres sont la graine ; l'étape de récurrence applique la composante \(\mathsf A\) de \(\operatorname{Upd}_{\varepsilon_N}\).

Adem\'as,
\begin{equation}
\label{eq:pdfv2-terminal-recuperacion-compatible}
 \operatorname{Rec}_{13;k,N}(h_{\leq N})
 =\operatorname{tr}_{13,N+1,N}
   \bigl(\operatorname{Rec}_{13;k,N+1}(h)\bigr),
\end{equation}
où \(\operatorname{tr}_{13,N+1,N}\) supprime seulement la dernière actualisation et restreint son cône de queues. En conséquence, un élément de \(\mathfrak T_{k,\infty}(\widehat x_k)\) détermine un système compatible des treize coordonnées à toute profondeur.

\section{Terminalité multivaluée, naturalité et limite compatible du continu}

\begin{theorem}[Terminal multivalué, naturalité et limite du continu]
\label{thm:pdfv2-terminal-conjunto}
Pour tout \(\widehat x_k\in\widehat{\mathcal X}^{\rm tot}_k\), les règles \eqref{eq:pdfv2-terminal-historias-finita}--\eqref{eq:pdfv2-terminal-recuperacion-compatible} construisent, à partir de l'unique continu discret :
\begin{enumerate}
 \item une multisection terminale finie non vide à chaque horizon ;
 \item une évaluation qui conserve toutes les histoires et multiplicités ;
 \item la connexion nonadique comme correspondance munie de témoins ;
 \item l'unique identité de naturalité \eqref{eq:pdfv2-terminal-naturalidad-unica} ;
 \item un terminal inverse non vide et une correspondance nonadique à la limite ;
 \item la reconstruction compatible de la fibre, des extensions, des relations, des nombres, de l'orientation, de la retenue, de la mémoire, de la frontière, du parcours, de la multisection, de la survie, du terminal et du lecteur.
\end{enumerate}
Les cinq caractéristiques intrinsèques restent corrélées au sein de chaque registre et de chaque témoin d'extension.
\end{theorem}

\begin{proof}
La non-vacuité de \(\mathscr H_{k,N}(\widehat x_k)\) est non vide par \cref{prop:continuo-comun-arbol-finito-no-vacio}, et sa troncature est surjective par \cref{lem:continuo-comun-terminal-sobreyectivo}. La règle \eqref{eq:pdfv2-terminal-actualizacion-unica} applique les identités APP, TRIT et TPK sur chaque témoin \(\varepsilon_j\) ; par récurrence sur \(N-k\), elle produit tous les champs de \eqref{eq:pdfv2-terminal-registro-trece} et conserve leurs relations.

Prendre l'ensemble complet dans \eqref{eq:pdfv2-terminal-multivaluado}, ou la somme complète dans \eqref{eq:pdfv2-terminal-evaluacion-todas}, démontre que le terminal ne contient pas de sélecteur. La version pondérée est isométrique par l'égalité projective \eqref{eq:pdfv2-terminal-isometria}.

Le témoin \(z=(\widehat x_k,\boldsymbol\omega)\) porte l'action de phase de chacune de ses neuf émissions ; par \cref{prop:continuo-comun-gamma9-retorno-memoria}, il définit les deux applications de \eqref{eq:pdfv2-terminal-span-gamma9} sans transformer la relation en fonction ni assimiler phase et état. Les formules \eqref{eq:pdfv2-terminal-gamma9-s-horizonte}--\eqref{eq:pdfv2-terminal-gamma9-enlace-registros} prouvent les deux carrés d'horizon et la même construction sur les registres. Sommer l'espace source de chaque témoin donne \eqref{eq:pdfv2-terminal-linealizacion-span} sans effacer les multiplicités.

La partition \eqref{eq:pdfv2-terminal-particion-historias} est disjointe puisque toute histoire possède un préfixe déterminé de longueur \(\ell\). Concaténer chaque préfixe avec toutes ses continuations reconstruit une fois et une seule chaque histoire de \(\mathscr H_{k,N}(\widehat x_k)\). Cela prouve directement \eqref{eq:pdfv2-terminal-naturalidad-unica} pour le registre entier ; cinq démonstrations disjointes ne sont pas nécessaires.

La surjectivité de \eqref{eq:pdfv2-terminal-transicion-horizonte} et la finitude de chaque niveau donnent la non-vacuité de \eqref{eq:pdfv2-terminal-limite-inverso}. Les liaisons explicites \eqref{eq:pdfv2-terminal-gamma9-enlace-apice} et \eqref{eq:pdfv2-terminal-gamma9-enlace-registros}, avec leurs carrés, construisent les limites \eqref{eq:pdfv2-terminal-gamma9-apice-limite-raw} et \eqref{eq:pdfv2-terminal-gamma9-apice-registros-limite} ; de celles-ci découle \eqref{eq:pdfv2-terminal-span-limite}. Enfin, \eqref{eq:pdfv2-terminal-recuperacion-trece} exécute à nouveau les actualisations successives, et \eqref{eq:pdfv2-terminal-recuperacion-compatible} démontre que ces lectures passent conjointement à la limite.
\end{proof}

\section{Obstructions à la terminalité de l'état enrichi : perte de fibre, d'orientation, de mémoire ou de lecteur}

\begin{proposition}[Falsificateurs du terminal conjoint]
\label{prop:pdfv2-terminal-falsadores}
Chacune des opérations suivantes remplace l'objet construit et bloque la conclusion du \cref{thm:pdfv2-terminal-conjunto} :
\begin{enumerate}
 \item envoyer un cylindre sur une seule histoire alors qu'il en admet plusieurs ;
 \item choisir une queue du champ terminal \(\mathsf T\) et supprimer les autres ;
 \item remplacer \(\Gamma_9\) par une fonction sans démontrer l'unicité de chaque cible et de son témoin ;
 \item linéariser une correspondance en ignorant les multiplicités ;
 \item identifier le retour de phase au retour d'état ;
 \item conserver l'entrée intacte comme première coordonnée et utiliser sa récupérabilité comme une génération supposée ;
 \item omettre l'un quelconque des treize champs de \eqref{eq:pdfv2-terminal-registro-trece} ;
 \item permettre à deux caractéristiques intrinsèques d'utiliser des histoires, retenues, frontières ou terminaux différents ;
 \item remplacer \eqref{eq:pdfv2-terminal-naturalidad-unica} par cinq affirmations indépendantes ;
 \item affirmer la compatibilité opératorielle de \(\Gamma_9\) sans conserver l'espace des témoins et leurs multiplicités ;
 \item intervertir sans preuve la formation des images et le passage à la limite ;
 \item présenter une liste de treize noms sans les actions \eqref{eq:pdfv2-terminal-actualizacion-unica} qui les produisent.
 \item omettre la télescopie polaire, la ligne des pertes, la naturalité adjointe, le module de toutes les continuations ou la réduction Smith--Bockstein établis dans \cref{prop:pdfv2-terminal-cinco-desarrollos-integrales} ;
 \item transformer l'un quelconque de ces cinq développements en branche autonome ou utiliser une application classique ultérieure pour sélectionner son domaine.
\end{enumerate}
\end{proposition}

\begin{proof}
Les cinq premiers actes contractent la multisection ou la correspondance. Les quatre suivants rompent le support et la naturalité communs. Les trois derniers suppriment la preuve générative ou modifient le passage à la limite. Dans chaque cas, l'objet résultant cesse de satisfaire au moins l'une de \eqref{eq:pdfv2-terminal-multivaluado}, \eqref{eq:pdfv2-terminal-span-terminal}, \eqref{eq:pdfv2-terminal-naturalidad-unica} ou \eqref{eq:pdfv2-terminal-recuperacion-compatible} ; la conclusion n'est donc pas transférable.
\end{proof}

\begin{statusbox}
L'organisation du terminal complet, la naturalité unique et le passage à la limite sont établis conjointement. La multisection survivante, les retenues, le registre, la frontière, le parcours et la connexion nonadique sont leurs antécédents structurels.\quad
La linéarisation sur les modules libres est démontrée par \eqref{eq:pdfv2-terminal-linealizacion-span} et ses carrés de Beck--Chevalley. Une promotion ultérieure à un opérateur borné sur une complétion analytique exige en outre de déclarer cette norme et de démontrer la borne ; la correspondance \eqref{eq:pdfv2-terminal-span-limite} ne dépend pas de cette interface.
\end{statusbox}
